\label{sec:conclusion_body}

Public Q\&A platforms codify routine problem-solving and
distribute spillovers across future learners. ChatGPT sharply
reduced the marginal cost of obtaining the same service
privately, and the recent literature
\citep{burtch2024chatgpt,delriochanona2024large,xue2026chatgpt}
documents the resulting aggregate contraction together with an
apparent improvement in surviving question quality---a
combination compatible with a benign reading under which only
the low-value end is pruned. We test that reading at the
within-tag value margin.

Using a triple difference on 8.4 million questions (100 top
tags, 261 weeks, seven question types), the central finding
rejects the \emph{strict} benign-pruning reading: a single
pre-specified shape test returns a negative average displacement
across the value distribution ($-0.167$, $p<10^{-4}$) that is
roughly constant across the low-to-mid crowd-validated range and
fades only at the rare elite tail, so the channel is not confined
to score-zero posts (the exclusive bins $0$ through $4$ are each
individually significant and Bonferroni-robust; a conservative
Romano--Wolf correction over the nested cumulative thresholds
places the gradient at the $10\%$ joint level). The volume
baseline is $-0.108$ ($p=0.021$), strengthening to $-0.140$
($p=0.005$) once the a-priori boundary category
advanced-architecture is excluded, and negative and significant
in all 32 cells of a specification curve. Because exact pre-trend
tests reject, we read the estimate as a conditional post-ChatGPT
association, not a natural experiment: its credibility rests on a
sign-reversing placebo distribution, monotone event-study growth,
a pre-trend-matched design ($-0.129$), and a HonestDiD breakdown
at $\bar M=0.75$. Re-estimating on the full 100-tag panel
extended through 2025 strengthens the effect to $-0.172$
($p=0.002$), deepening to $-0.68$ with no reversion---most
consistent with sustained diffusion rather than a transient
episode.

We turn the result into a \emph{value-calibrated early-warning
system}: out of sample, an activity-decline signal detects where
high-value erosion is largest (AUC $0.98$ versus $0.57$ for a
value-weighted detector) and the value-margin estimates
calibrate that alert into the high-value posts at risk, with a
persistent composition flag ($\rho=0.90$) classifying
intervention type. We do not document an effect on innovation
\emph{outcomes}, contributor entry, or productivity, and we
derive no welfare conclusion; the Q\&A margin is an innovation
\emph{input}, and the private productivity gains of LLM adoption
are not estimated. The conceptual implication is that generative
AI changes the \emph{value composition}, not only the quantity,
of the public knowledge inputs that fuel open innovation---the
dimension that innovation policy and management research should
follow. The entire pipeline is reproducible from public,
zero-budget data; the full robustness, mechanism, and validation
detail is in the replication package.
